% ============================================================
\chapter{Modul 8: Module und Pakete}
% ============================================================

\section{Lektion 22: Bibliotheken nutzen}

\begin{lstlisting}
# Standardbibliothek -- immer verfuegbar
import math
import time
import random
import os

# Mathematik für Roboterkinematik
winkel_rad = math.radians(90)      # Grad -> Bogenmaß
winkel_deg = math.degrees(1.5708)  # Bogenmaß -> Grad
print(f"sin(90 Grad) = {math.sin(winkel_rad):.4f}")

# Euklidischer Abstand (Wegpunkte)
def abstand(p1, p2):
    dx = p2[0] - p1[0]
    dy = p2[1] - p1[1]
    return math.sqrt(dx**2 + dy**2)

aktuell = (0.0, 0.0)
ziel    = (3.0, 4.0)
print(f"Entfernung zum Ziel: {abstand(aktuell, ziel)} m")

# Zeit und Timing
start = time.time()
time.sleep(0.1)   # 100ms warten (Sensor-Abtastrate)
dauer = time.time() - start
print(f"Messzeit: {dauer*1000:.1f} ms")
\end{lstlisting}

\section{Lektion 23: Eigene Module erstellen}

\begin{lstlisting}
# Datei: roboter_utils.py
"""
Hilfsfunktionen fuer den MobileRobot-1
"""

import math

RADABSTAND = 0.15   # Meter, Konstante

def normiere_winkel(winkel_deg):
    """Winkel auf -180 bis +180 Grad normieren."""
    while winkel_deg > 180:
        winkel_deg -= 360
    while winkel_deg < -180:
        winkel_deg += 360
    return winkel_deg

def v_zu_pwm(v, v_max=0.5, pwm_max=255):
    """Lineare Konvertierung: m/s zu PWM-Wert."""
    return int((v / v_max) * pwm_max)

def kinematik_vorwaerts(v_links, v_rechts, dt=0.1):
    """Differentialantrieb-Kinematik: Geschwindigkeit aus Radwerten."""
    v     = (v_links + v_rechts) / 2
    omega = (v_rechts - v_links) / RADABSTAND
    return v, omega
\end{lstlisting}

\begin{lstlisting}
# Datei: main.py -- eigenes Modul importieren
from roboter_utils import v_zu_pwm, kinematik_vorwaerts

pwm = v_zu_pwm(0.3)
print(f"PWM-Wert: {pwm}")

v, omega = kinematik_vorwaerts(0.30, 0.35)
print(f"v={v:.3f} m/s, omega={omega:.3f} rad/s")
\end{lstlisting}

\begin{merksatz}
Modulname = Dateiname ohne \texttt{.py}. \texttt{import roboter\_utils} lädt alles. \texttt{from roboter\_utils import v\_zu\_pwm} lädt nur eine Funktion. Eigene Module gehören in dasselbe Verzeichnis wie \texttt{main.py}.
\end{merksatz}

\begin{ros}
In ROS2 ist jedes Package ein Python-Modul (oder C++-Paket). Die Struktur ist formalisiert: \texttt{setup.py}, \texttt{package.xml}, \texttt{resource/}. Aber das Prinzip ist dasselbe -- wiederverwendbarer, importierbarer Code.
\end{ros}

\begin{uebung}[title=Projektaufgabe Modul 8: Roboter-Modul]
Erstelle eine Datei \texttt{mobile\_robot.py} mit Konstanten (Radabstand, Max-Geschwindigkeit) und mindestens 5 Funktionen (Kinematik, Normierung, Einheitenkonvertierung). Importiere das Modul in \texttt{main.py} und demonstriere alle Funktionen.
\end{uebung}
